% ID: SEL-D0061-OA-E02
% ESTADO: cotejada. Segunda pasada visual y borrador compilado revisados.
% Fuente: Selectividad/sel_2010_matematicas_aplicadas/mat_apli/mat_apli_junio/mat_apli_junio/EXAMEN-ANDALUCÍA-5.pdf
% SHA-256: 5ea9fab18d8185843726e08994f39cf34f9b7dab32d86c310902a2f203dace2a
% Página: 1; opción A; ejercicio original 2.

Sea la función $f(x)=2x^2-\dfrac{1}{3}x^3$. Calcule:
\begin{enumerate}[label=\alph*)]
\item \textbf{(1 punto)} Los intervalos de crecimiento y decrecimiento.
\item \textbf{(1 punto)} Las coordenadas de sus extremos relativos.
\item \textbf{(0.5 puntos)} El punto de la gráfica en el que la pendiente de la recta tangente a dicha gráfica es 4.
\end{enumerate}
